% !TEX root = ../main.tex
\begin{ouverture}[transferts]{16}{Énergie potentielle et énergie mécanique}{Partie 3 · L'énergie : conversions et transferts}
{Énergie potentielle de pesanteur\enspace·\enspace Énergie mécanique\enspace·\enspace Conservation\enspace·\enspace Frottements et dissipation}
\objectif{Établir et utiliser l'expression de l'énergie potentielle de pesanteur au voisinage de la Terre.}
\objectif{Identifier des situations de conservation et de non-conservation de l'énergie mécanique.}
\objectif{Exploiter la conservation de l'énergie mécanique : chute libre, pendule sans frottement.}
\objectif{Utiliser la variation de l'énergie mécanique pour déterminer le travail des forces non conservatives ; faire un bilan avec un programme.}
\prerequis{Énergie cinétique, travail d'une force, théorème de l'énergie cinétique (chapitre 15).}
\prerequis{Travail du poids $W_{\mathrm{AB}}(\vec{P})=m\,g\,(z_\mathrm{A}-z_\mathrm{B})$ ; forces conservatives et non conservatives.}
\end{ouverture}
\essentiel
\essSection{1}{Énergie potentielle de pesanteur}
\begin{definition}[Énergie potentielle de pesanteur]
Le travail du poids ne dépend que des altitudes : $W_{\mathrm{AB}}(\vec{P})=m\,g\,z_\mathrm{A}-m\,g\,z_\mathrm{B}$. On définit l'énergie potentielle de pesanteur, au
voisinage de la Terre, avec un axe $z$ vertical orienté vers le haut :
\[E_\mathrm{p}=m\,g\,z\qquad\text{(en J), et alors}\qquad W_{\mathrm{AB}}(\vec{P})=-\Delta E_\mathrm{p}=E_\mathrm{p}(\mathrm{A})-E_\mathrm{p}(\mathrm{B}).\]
Elle dépend du choix de l'origine des altitudes ($E_\mathrm{p}=0$ en $z=0$), mais ses variations n'en dépendent pas. Un livre de \qty{1,0}{kg}
posé \qty{2,0}{m} au-dessus du sol a, par rapport au sol, une énergie potentielle de \qty{19,6}{J}.
\end{definition}
\essSection{2}{Énergie mécanique}
\begin{definition}[Énergie mécanique]
L'énergie mécanique d'un système est la somme de ses énergies cinétique et potentielle : $E_\mathrm{m}=E_\mathrm{c}+E_\mathrm{p}=\dfrac{1}{2}\,m\,v^2+m\,g\,z$.
\end{definition}
\begin{loi}[Conservation et non-conservation]
D'après le théorème de l'énergie cinétique, la variation de l'énergie mécanique est égale au travail des forces non
conservatives :
\[\Delta E_\mathrm{m}=E_\mathrm{m}(\mathrm{B})-E_\mathrm{m}(\mathrm{A})=W_{\mathrm{AB}}(\vec{F}_\text{non conservatives}).\]
Si seul le poids travaille (chute libre, pendule ou glissade sans frottement), l'énergie mécanique se \textbf{conserve}. Des
frottements la font \textbf{diminuer} (énergie dissipée, le plus souvent thermique) ; un moteur peut la faire augmenter.
\end{loi}
\schema{16-conservation}
\begin{propriete}[Exploiter la conservation]
Sans frottement, entre deux positions : $\dfrac{1}{2}\,m\,v_\mathrm{A}^2+m\,g\,z_\mathrm{A}=\dfrac{1}{2}\,m\,v_\mathrm{B}^2+m\,g\,z_\mathrm{B}$. La masse se simplifie : un objet lâché de \qty{10}{m}
arrive au sol à $\sqrt{2\,g\,h}=\qty{14,0}{m.s^{-1}}$. Pour un pendule, l'énergie passe sans cesse de la forme potentielle à la forme cinétique.
\end{propriete}
\par\penalty0
\schema{16-pendule}
\begin{python}
m, g = 0.60, 9.81                   # kg, N/kg : ballon lache de 2,00 m
z = [2.00, 1.80, 1.22, 0.23]        # m, toutes les 0,20 s
v = [0.00, 1.96, 3.92, 5.89]        # m/s
for i in range(4):                  # energies en J
    Ec, Ep = 0.5 * m * v[i]**2, m * g * z[i]
    print('Ec =', round(Ec, 2), '| Ep =', round(Ep, 2), '| Em =', round(Ec + Ep, 2))
\end{python}
\begin{console}
Ec = 0.0 | Ep = 11.77 | Em = 11.77
Ec = 1.15 | Ep = 10.59 | Em = 11.75
Ec = 4.61 | Ep = 7.18 | Em = 11.79
Ec = 10.41 | Ep = 1.35 | Em = 11.76
\end{console}
\noindent\capacite{Capacité numérique} Le bilan énergétique, calculé à partir des mesures, montre que l'énergie mécanique du ballon reste constante
aux erreurs de mesure près : l'action de l'air est négligeable sur cette chute.
\schema{16-frottements}
\begin{attention}
Toujours préciser l'origine choisie pour les altitudes, et la garder pendant tout l'exercice. L'énergie mécanique ne se conserve
que si les forces non conservatives (frottements, poussée d'un moteur, action d'un opérateur) ne travaillent pas.
\end{attention}
\begin{experience}[Bilan énergétique d'une chute ou d'un pendule]
Avec une vidéo et un logiciel de pointage, ou les capteurs d'un smartphone, on obtient positions et vitesses ; un tableur ou un
programme calcule $E_\mathrm{c}$, $E_\mathrm{p}$ et $E_\mathrm{m}$. On compare une balle lourde et une balle légère, ou un pendule au fil des oscillations.
\end{experience}
\begin{rediger}
Définir le système, le référentiel et l'origine des altitudes ; écrire l'énergie mécanique aux deux positions ; justifier la
conservation (seul le poids travaille) ou écrire $\Delta E_\mathrm{m}=W(\vec{f})$ avant de calculer.
\end{rediger}
\exercices

\rubrique{Découverte}
\begin{blocLivre}{ex-16-1}
\begin{exercice}[1]{Calculer des énergies potentielles}
\begin{donnees}
$g=\qty{9,81}{N.kg^{-1}}$.
\end{donnees}
\begin{questions}
\item Calculer l'énergie potentielle d'un livre de \qty{0,80}{kg} posé sur une étagère à \qty{2,0}{m} du sol, l'origine étant prise au sol.
\item Que vaut l'énergie potentielle du livre posé au sol si l'on prend l'origine des altitudes au niveau de l'étagère ?
\item Calculer l'énergie potentielle d'une alpiniste de \qty{70}{kg} au sommet du mont Blanc (\qty{4808}{m}), l'origine étant au niveau de la mer.
\item Calculer sa variation d'énergie potentielle depuis Chamonix (\qty{1035}{m}). Dépend-elle de l'origine choisie ?
\end{questions}
\end{exercice}
\end{blocLivre}
\begin{blocLivre}{ex-16-2}
\begin{exercice}[1]{Une chute sans frottement}
On lâche une balle de \qty{0,20}{kg} sans vitesse d'une hauteur de \qty{10}{m}. On néglige l'action de l'air ; l'origine des altitudes est au sol.
\begin{donnees}
$g=\qty{9,81}{N.kg^{-1}}$.
\end{donnees}
\begin{questions}
\item Calculer l'énergie mécanique de la balle au moment où on la lâche.
\item Pourquoi l'énergie mécanique se conserve-t-elle pendant la chute ?
\item Calculer les énergies potentielle et cinétique de la balle à \qty{4,0}{m} du sol, puis sa vitesse.
\item Calculer la vitesse de la balle à l'arrivée au sol.
\end{questions}
\end{exercice}
\end{blocLivre}
\begin{blocLivre}{ex-16-3}
\begin{exercice}[1]{Lire des courbes d'énergie}
Une balle de \qty{0,10}{kg} est lancée verticalement vers le haut à \qty{8,0}{m.s^{-1}}, depuis \qty{1,0}{m} au-dessus du sol.
\schema{16-ex03}
\begin{questions}
\item Associer chaque courbe à l'énergie qu'elle représente. Justifier.
\item Lire les valeurs initiales des énergies et vérifier par le calcul ($g=\qty{9,81}{N.kg^{-1}}$).
\item À quelle date la balle atteint-elle son point le plus haut ? Calculer son altitude.
\item L'action de l'air est-elle négligeable ? Justifier.
\end{questions}
\end{exercice}
\end{blocLivre}
\begin{blocLivre}{ex-16-4}
\begin{exercice}[1]{Le pendule simple}
\schema{16-ex04}
\begin{donnees}
$g=\qty{9,81}{N.kg^{-1}}$.
\end{donnees}
\begin{questions}
\item Montrer que la bille descend d'une hauteur $h=L\,(1-\cos\ang{30})$ entre A et B, et la calculer.
\item Quelles forces s'exercent sur la bille ? Pourquoi la tension du fil ne travaille-t-elle pas ?
\item Calculer la vitesse de la bille au passage en B.
\item Calculer cette vitesse si on lâche la bille à \ang{60} de la verticale.
\end{questions}
\end{exercice}
\end{blocLivre}
\begin{blocLivre}{ex-16-5}
\begin{exercice}[1]{Conservation ou non ?}
Pour chaque situation, dire si l'énergie mécanique du système se conserve, augmente ou diminue, et justifier.
\begin{questions}
\item Une pomme tombe d'un arbre ; on néglige l'action de l'air.
\item Un parachutiste descend à vitesse constante.
\item Une skieuse remonte la piste, tirée à vitesse constante par un téléski.
\item Une bille oscille au bout d'un fil ; l'amplitude diminue peu à peu.
\end{questions}
\end{exercice}
\end{blocLivre}
\rubrique{Entraînement}
\begin{blocLivre}{ex-16-6}
\begin{exercice}[2]{La balle de ping-pong, en Python}
On lâche une balle de ping-pong de \qty{2,7}{g} de \qty{2,0}{m} ; une vidéo donne son altitude et sa vitesse toutes les \qty{0,10}{s}.
\begin{python}
m, g = 0.0027, 9.81
z = [2.000, 1.951, 1.805, 1.566, 1.239, 0.830, 0.347]   # m
v = [0, 0.98, 1.93, 2.84, 3.69, 4.47, 5.17]             # m/s
for i in range(7):
    Em = 0.5 * m * v[i]**2 + m * g * z[i]
    print(round(0.1 * i, 1), 's :', round(Em * 1000, 1), 'mJ')
\end{python}
\begin{console}
0.0 s : 53.0 mJ
0.1 s : 53.0 mJ
0.2 s : 52.8 mJ
0.3 s : 52.4 mJ
0.4 s : 51.2 mJ
0.5 s : 49.0 mJ
0.6 s : 45.3 mJ
\end{console}
\begin{questions}
\item \capacite{Capacité numérique}Quelle grandeur le programme calcule-t-il ? Quelle est l'origine des altitudes ?
\item Comment évolue l'énergie mécanique ? Interpréter.
\item Calculer le travail des frottements de l'air entre 0 et \qty{0,6}{s}.
\item Sans frottement, la vitesse à \qty{0,6}{s} vaudrait $g\,t$. Comparer à la vitesse mesurée.
\end{questions}
\end{exercice}
\end{blocLivre}
\begin{blocLivre}{ex-16-7}
\begin{exercice}[2]{Les montagnes russes}
\schema{16-ex07}
\begin{donnees}
$g=\qty{9,81}{N.kg^{-1}}$.
\end{donnees}
\begin{questions}
\item Calculer l'énergie mécanique du wagon en A, l'origine des altitudes étant au sol.
\item En négligeant les frottements, calculer la vitesse du wagon en B, puis en C.
\item La vitesse mesurée en B vaut \qty{24}{m.s^{-1}}. Calculer l'énergie dissipée entre A et B.
\item Quelle fraction de l'énergie mécanique initiale cela représente-t-il ?
\end{questions}
\end{exercice}
\end{blocLivre}
\begin{blocLivre}{ex-16-8}
\begin{exercice}[2]{Les rebonds d'une balle}
Une balle de tennis de \qty{58}{g}, lâchée de \qty{1,20}{m}, remonte à \qty{0,80}{m} après le premier rebond.
\begin{donnees}
$g=\qty{9,81}{N.kg^{-1}}$.
\end{donnees}
\begin{questions}
\item Comparer les énergies mécaniques au point de départ et au sommet du premier rebond.
\item Calculer l'énergie perdue lors du rebond. Sous quelle forme est-elle dissipée ?
\item Quelle fraction de son énergie mécanique la balle conserve-t-elle à chaque rebond ?
\item En supposant cette fraction constante, à quelle hauteur la balle remonte-t-elle après cinq rebonds ?
\end{questions}
\end{exercice}
\end{blocLivre}
\begin{blocLivre}{ex-16-9}
\begin{exercice}[2]{Une balle lancée vers le haut}
Une balle de \qty{150}{g} est lancée verticalement vers le haut à \qty{12}{m.s^{-1}}.
\begin{donnees}
$g=\qty{9,81}{N.kg^{-1}}$.
\end{donnees}
\begin{questions}
\item En négligeant l'action de l'air, calculer la hauteur maximale atteinte au-dessus du point de lancer.
\item La balle ne monte en réalité qu'à \qty{6,5}{m}. Calculer le travail des frottements de l'air pendant la montée.
\item Sans frottement, quelle serait la vitesse de la balle en repassant par le point de lancer ?
\item Avec frottement, sera-t-elle plus grande ou plus petite ? Justifier.
\end{questions}
\end{exercice}
\end{blocLivre}
\begin{blocLivre}{ex-16-10}
\begin{exercice}[2]{Le téléski}
Un téléski tire à vitesse constante une skieuse de \qty{75}{kg} sur une piste dont le dénivelé est de \qty{300}{m} ; la montée dure \qty{6,0}{min}.
\begin{donnees}
$g=\qty{9,81}{N.kg^{-1}}$.
\end{donnees}
\begin{questions}
\item Comment évoluent l'énergie cinétique, l'énergie potentielle et l'énergie mécanique de la skieuse ?
\item Calculer la variation de son énergie mécanique pendant la montée.
\item En négligeant les frottements, quel travail le téléski fournit-il ? Calculer la puissance correspondante.
\item En réalité, le travail du téléski est plus grand. Pourquoi ?
\end{questions}
\end{exercice}
\end{blocLivre}
\begin{blocLivre}{ex-16-11}
\begin{exercice}[2]{Le pendule réel}
\schema{16-ex11}
\begin{donnees}
$g=\qty{9,81}{N.kg^{-1}}$.
\end{donnees}
\begin{questions}
\item Identifier les courbes 1, 2 et 3. Pourquoi les courbes 1 et 2 oscillent-elles deux fois par période du pendule ?
\item Vérifier par le calcul la valeur initiale de l'énergie mécanique.
\item Après \qty{20}{s}, l'amplitude n'est plus que de \ang{30}. Calculer l'énergie mécanique, puis l'énergie dissipée.
\item Calculer la puissance moyenne dissipée par les frottements.
\end{questions}
\end{exercice}
\end{blocLivre}
\rubrique{Objectif DS}
\begin{blocLivre}{ex-16-12}
\begin{exercice}[3]{La centrale hydroélectrique}
\begin{docu}{1}{Une chute d'eau}
Dans une centrale de montagne, l'eau d'un lac de retenue descend dans une conduite forcée de \qty{300}{m} de dénivelé jusqu'aux turbines,
avec un débit de \qty{50}{m^3} par seconde. L'ensemble turbine-alternateur a un rendement de \qty{90}{\percent}.
\end{docu}
\begin{docu}{2}{Consommation}
Un foyer français consomme en moyenne \qty{4700}{kWh} d'électricité par an.
\end{docu}
\begin{donnees}
$g=\qty{9,81}{N.kg^{-1}}$ ; masse volumique de l'eau : \qty{1000}{kg.m^{-3}}.
\end{donnees}
\begin{questions}
\item Calculer l'énergie potentielle perdue par \qty{1}{m^3} d'eau en descendant de \qty{300}{m}.
\item En déduire la puissance cédée par l'eau aux turbines, puis la puissance électrique produite.
\item Calculer l'énergie électrique produite en un jour, en \unit{MWh}.
\item \initiative Combien de foyers cette centrale pourrait-elle alimenter, en fonctionnant ainsi toute l'année ?
\end{questions}
\end{exercice}
\end{blocLivre}
\begin{blocLivre}{ex-16-13}
\begin{exercice}[3]{Le lancer du poids}
\begin{docu}{1}{Le lancer}
Un athlète lâche le poids de \qty{7,26}{kg} à \qty{2,1}{m} du sol, avec une vitesse de \qty{13}{m.s^{-1}} inclinée de \ang{40} au-dessus de l'horizontale. On
néglige l'action de l'air : la composante horizontale de la vitesse reste alors constante.
\end{docu}
\begin{donnees}
$g=\qty{9,81}{N.kg^{-1}}$ ; origine des altitudes au sol.
\end{donnees}
\begin{questions}
\item Calculer l'énergie mécanique du poids au moment du lâcher.
\item Calculer la vitesse du poids à l'arrivée au sol.
\item Au sommet de la trajectoire, la vitesse est horizontale et vaut $13\times\cos\ang{40}$. Calculer l'énergie cinétique au sommet.
\item \initiative En déduire l'altitude du sommet de la trajectoire.
\end{questions}
\end{exercice}
\end{blocLivre}
\begin{blocLivre}{ex-16-14}
\begin{exercice}[3]{La descente à vélo}
\begin{docu}{1}{Une longue descente}
Une cycliste (\qty{80}{kg} avec son vélo) descend sans pédaler une route rectiligne de pente \qty{10}{\percent} : la route descend de \qty{10}{m} pour \qty{100}{m}
parcourus à l'horizontale, ce qui correspond à un angle dont le sinus vaut \num{0,0995}. Sa vitesse se stabilise à \qty{15}{m.s^{-1}}.
\end{docu}
\begin{donnees}
$g=\qty{9,81}{N.kg^{-1}}$.
\end{donnees}
\begin{questions}
\item À vitesse constante, comment évoluent les énergies cinétique, potentielle et mécanique ?
\item De combien la cycliste descend-elle en une seconde ?
\item \initiative En déduire la puissance dissipée par les frottements (air et roulement).
\item Sans frottement, quelle serait sa vitesse après \qty{2}{km} de descente, si elle partait à \qty{15}{m.s^{-1}} ? Commenter.
\end{questions}
\end{exercice}
\end{blocLivre}
\begin{blocLivre}{ex-16-15}
\begin{exercice}[3]{Le bilan d'énergie des rebonds, en Python}
\begin{docu}{1}{Des mesures}
Une balle lâchée de \qty{1,50}{m} remonte successivement à \qty{1,10}{m}, \qty{0,81}{m}, \qty{0,59}{m} puis \qty{0,43}{m}.
\end{docu}
\begin{docu}{2}{Le programme}
\begin{python}
h = [1.50, 1.10, 0.81, 0.59, 0.43]          # m : hauteurs atteintes
for i in range(4):
    print('rebond', i + 1, ':', round(h[i + 1] / h[i], 2))
k = 0.73
hauteur, n = h[-1], 4
while hauteur > 0.05:
    hauteur = hauteur * k
    n = n + 1
print(n, 'rebonds')
\end{python}
\begin{console}
rebond 1 : 0.73
rebond 2 : 0.74
rebond 3 : 0.73
rebond 4 : 0.73
11 rebonds
\end{console}
\end{docu}
\begin{questions}
\item \capacite{Capacité numérique}Pourquoi le rapport de deux hauteurs successives est-il aussi le rapport des énergies mécaniques ?
\item Quelle fraction de son énergie mécanique la balle perd-elle à chaque rebond ?
\item Que calcule la boucle \code{while} ? Interpréter le résultat affiché.
\item \initiative La balle a une masse de \qty{58}{g}. Calculer l'énergie dissipée lors du premier rebond ($g=\qty{9,81}{N.kg^{-1}}$).
\end{questions}
\end{exercice}
\end{blocLivre}
